
\subsection*{1.5}
\phantomsection
\label{sga3:VIIB:1.5}
\addcontentsline{toc}{subsection}{1.5. The flat topology}

Let \(k\) be a pseudocompact ring.

\medskip
\noindent\textbf{Definition.}
\phantomsection
\label{sga3:VIIB:1.5.definition}
We shall say that a family of morphisms
\[
  f_i:X_i\longrightarrow X
\]
in \(\operatorname{Vaf}/k\) is a \emph{surjective topologically flat
family} if the morphism
\[
  \coprod_i X_i\longrightarrow X
\]
induced by the \(f_i\) is surjective and topologically flat. This
means that every \(f_i\) is topologically flat and that every point of
\(X\) lies in the image of at least one \(X_i\).

It follows from \SGARefLink{sga3:VIIB:1.3.3}{1.3.3} that surjective topologically flat families
define a pretopology on \(\operatorname{Vaf}/k\) \textup{(IV,
\SGARefLink{sga3:iv:definition:4:2:5}{4.2.5})}. The corresponding topology will be called the \emph{flat
topology} on \(\operatorname{Vaf}/k\).

By IV, \SGARefLink{sga3:iv:corollary:4:3:5}{4.3.5}, a functor
\[
  F:(\operatorname{Vaf}/k)^{\mathrm{op}}\longrightarrow
  (\mathbf{Sets})
\]
is a sheaf for the flat topology if and only if \(F\) carries every
direct sum to a direct product and the sequence
\begin{center}
  \phantomsection
  \label{sga3:VIIB:1.5:flat-sheaf-sequence}
  \(
F(Y)\xrightarrow{\ F(f)\ }F(X)
\mathrel{\substack{\xrightarrow{\ F(\operatorname{pr}_1)\ }\\[-.7ex]
                   \xrightarrow[\ F(\operatorname{pr}_2)\ ]{}}}
F(X\times_YX)
\)

\end{center}
is exact for every surjective topologically flat morphism
\(f:X\to Y\).

By IV, \SGARefLink{sga3:iv:subsection:4:5}{4.5},\footnote{\label{sga3:VIIB:nde:69}The editors modified the
original in what follows. In particular, they replaced the statement:
``if \(d_0,d_1:X_1\rightrightarrows X\) is an equivalence relation
such that \(d_1\) is topologically flat, formation of the quotient
commutes with \(h\),'' by the theorem below.}
\SGARefLink{sga3:VIIB:1.3.1}{Proposition~1.3.1} implies that the flat topology is coarser than the
canonical topology. In other words, for every object \(X\) of
\(\operatorname{Vaf}/k\), the functor
\[
  h_X:T\longmapsto
  \operatorname{Hom}_{\operatorname{Vaf}/k}(T,X)
\]
is a sheaf for the flat topology. In what follows, as usual
\textup{(cf.\ Exposé I)}, we identify \(X\) with \(h_X\).

By IV, \SGARefLink{sga3:iv:lemma:4:6:5}{4.6.5}, \SGARefLink{sga3:VIIB:1.4.thm}{Theorem~1.4} may be reformulated as follows.

\medskip
\noindent\textbf{Theorem.}
\phantomsection
\label{sga3:VIIB:1.5.theorem}
\textit{Let \(k\) be a pseudocompact ring, let
\(d_0,d_1:X_1\rightrightarrows X\) be an equivalence pair in
\(\operatorname{Vaf}/k\), and let \(X/X_1\) be the quotient formal
variety \textup{(that is,
\(\operatorname{Coker}(d_0,d_1)\), cf.\ \SGARefLink{sga3:VIIB:1.2}{1.2})}. If \(d_1\) is
topologically flat, then \(X/X_1\) represents the quotient sheaf for
the flat topology.}
